\documentclass[11pt,a4paper]{article}

% ------------------------------------------------------------------------------
% Packages
% ------------------------------------------------------------------------------
\usepackage[utf8]{inputenc}
\usepackage[T1]{fontenc}
\usepackage[french]{babel}
\usepackage{lmodern}
\usepackage[a4paper,margin=2.2cm]{geometry}
\usepackage{xcolor}
\usepackage{array}
\usepackage{multirow}
\usepackage{float}
\usepackage{caption}
\usepackage{listings}
\usepackage{tikz}
\usetikzlibrary{babel}
\usepackage{hyperref}

\hypersetup{colorlinks=true, linkcolor=bleu, urlcolor=bleu}

% ------------------------------------------------------------------------------
% Couleurs
% ------------------------------------------------------------------------------
\definecolor{bleu}{HTML}{1E3A8A}
\definecolor{rougenoeud}{HTML}{DC2626}
\definecolor{noirnoeud}{HTML}{1F2937}
\definecolor{vert}{HTML}{16A34A}
\definecolor{gris}{HTML}{64748B}
\definecolor{codefond}{HTML}{F8FAFC}
\definecolor{codecadre}{HTML}{CBD5E1}

% ------------------------------------------------------------------------------
% Style du code C++
% ------------------------------------------------------------------------------
\lstset{
  language=C++,
  basicstyle=\ttfamily\footnotesize,
  keywordstyle=\color{bleu}\bfseries,
  commentstyle=\color{gris}\itshape,
  stringstyle=\color{vert},
  backgroundcolor=\color{codefond},
  frame=single,
  rulecolor=\color{codecadre},
  breaklines=true,
  columns=fullflexible,
  keepspaces=true,
  showstringspaces=false,
  numbers=left,
  numberstyle=\tiny\color{gris},
  captionpos=b,
}
\renewcommand{\lstlistingname}{Code}

% ------------------------------------------------------------------------------
% Noeuds de l'arbre
% ------------------------------------------------------------------------------
\tikzset{
  rn/.style={circle, minimum size=8mm, inner sep=0pt, font=\small\bfseries, text=white},
  rouge/.style={rn, fill=rougenoeud},
  noir/.style={rn, fill=noirnoeud},
  edge from parent/.style={draw=black!45, thick},
}

\setlength{\parindent}{0pt}
\setlength{\parskip}{5pt}
\renewcommand{\arraystretch}{1.2}
\newcommand{\entete}[1]{\textbf{\begin{tabular}[c]{@{}c@{}}#1\end{tabular}}}

\title{\textbf{Arbre rouge-noir}\\[4pt]
       \large Insertion, recherche et suppression, comparaison avec \texttt{std::map}
       et \texttt{std::unordered\_map}}
\author{Zakariae Bakkari\\
        \small Programmation Objet Avancée et Design Patterns\\
        \small\url{https://github.com/zakariae-bakkari/cpp_ilisi2}}
\date{11 novembre 2026}

\begin{document}

\maketitle

% ==============================================================================
\section{Objectif}
% ==============================================================================

Écrire un \textbf{arbre rouge-noir} en C++, sans classe (une \texttt{struct Noeud} et des
fonctions), y insérer des valeurs, puis mesurer le temps d'insertion, de recherche et la
hauteur de l'arbre. Ensuite, le comparer à \texttt{std::map} (l'arbre rouge-noir de la STL)
et à \texttt{std::unordered\_map} (une table de hachage).

\textbf{Code source :}
\begin{itemize}\footnotesize
  \item arbre rouge-noir, dessin coloré et temps (option \textbf{15} du menu
        \texttt{./run.sh}, ou \texttt{make run\_14}) :\\
        \url{https://github.com/zakariae-bakkari/cpp_ilisi2/blob/main/iterators/src/14_arbre_rouge_noir.cpp}
  \item comparaison avec \texttt{std::map} et \texttt{std::unordered\_map} (option
        \textbf{16}, ou \texttt{make run\_15}) :\\
        \url{https://github.com/zakariae-bakkari/cpp_ilisi2/blob/main/iterators/src/15_arbre_rouge_noir_vs_map.cpp}
\end{itemize}

% ==============================================================================
\section{Les règles}
% ==============================================================================

Un arbre rouge-noir est un arbre binaire de recherche dont chaque nœud porte une couleur.
Les feuilles vides (\texttt{nil}, ici \texttt{nullptr}) comptent comme noires.

\begin{table}[H]
\centering
\begin{tabular}{|c|l|}
\hline
\textbf{Règle} & \textbf{Énoncé} \\
\hline
1 & Chaque nœud est rouge ou noir. \\
\hline
2 & La racine est noire. \\
\hline
3 & Les feuilles \texttt{nil} sont noires. \\
\hline
4 & Un nœud rouge n'a pas d'enfant rouge. \\
\hline
5 & Tous les chemins d'un nœud vers les \texttt{nil} ont le même nombre de nœuds noirs. \\
\hline
\end{tabular}
\caption{Les 5 règles d'un arbre rouge-noir.}
\end{table}

\textbf{Pourquoi l'arbre reste équilibré.} Avec $h_n$ noirs sur chaque chemin (règle 5), le
plus court chemin n'a que des noirs ($h_n$ nœuds), et le plus long alterne noir et rouge
(au plus $2h_n$ nœuds, règle 4). Le plus long chemin fait donc au plus deux fois le plus
court, et la hauteur reste sous $2\log_2(n+1)$.

% ==============================================================================
\section{Insertion}
% ==============================================================================

On insère d'abord comme dans un arbre binaire de recherche. Le nouveau nœud N est
\textbf{rouge} : il n'ajoute pas de noir, donc la règle 5 reste vraie. Seule la règle 4
peut casser (si le parent P est rouge). On regarde alors l'oncle U, le frère de P :

\begin{table}[H]
\centering
\begin{tabular}{|c|p{4.6cm}|p{8.6cm}|}
\hline
\textbf{Cas} & \textbf{Situation} & \textbf{Action} \\
\hline
A & oncle rouge & recoloration : P et U noirs, grand-père G rouge, puis on recommence en G \\
\hline
B & oncle noir, forme en ligne & rotation en G, puis échange des couleurs de P et G \\
\hline
C & oncle noir, forme en zig-zag & rotation en P pour obtenir une ligne, puis cas B \\
\hline
\end{tabular}
\caption{Les trois cas de réparation (et leurs miroirs quand P est à droite de G).}
\end{table}

À la fin, la racine est remise en noir (règle 2).

\begin{lstlisting}[caption={Réparation après insertion (\texttt{14\_arbre\_rouge\_noir.cpp}).}]
void reparer(Noeud*& racine, Noeud* n) {
    while (estRouge(n->parent)) {
        Noeud* p = n->parent;
        Noeud* g = p->parent;              // existe : un parent rouge n'est pas la racine
        bool pGauche = (p == g->gauche);   // false = cas miroir
        Noeud* oncle = pGauche ? g->droite : g->gauche;

        if (estRouge(oncle)) {             // oncle rouge : recoloration
            p->rouge = oncle->rouge = false;
            g->rouge = true;
            n = g;                         // on remonte
            continue;
        }
        if (n == (pGauche ? p->droite : p->gauche)) {  // zig-zag -> ligne
            pGauche ? rotationGauche(racine, p) : rotationDroite(racine, p);
            n = p;
            p = n->parent;
        }
        pGauche ? rotationDroite(racine, g) : rotationGauche(racine, g);  // ligne
        p->rouge = false;
        g->rouge = true;
    }
    racine->rouge = false;                 // la racine est toujours noire
}
\end{lstlisting}

% ==============================================================================
\section{Exemple : 10 valeurs}
% ==============================================================================

On insère, dans cet ordre : 1, 45, 67, 89, 12, 34, 23, 78, 56, 90. L'arbre obtenu :

\begin{figure}[H]
\centering
\begin{tikzpicture}[
  level distance=13mm,
  level 1/.style={sibling distance=60mm},
  level 2/.style={sibling distance=30mm},
  level 3/.style={sibling distance=15mm}]
\node[noir] {45}
  child { node[rouge] {12}
    child { node[noir] {1} }
    child { node[noir] {34}
      child { node[rouge] {23} }
      child[missing] }
  }
  child { node[rouge] {78}
    child { node[noir] {67}
      child { node[rouge] {56} }
      child[missing] }
    child { node[noir] {89}
      child[missing]
      child { node[rouge] {90} } }
  };
\end{tikzpicture}
\caption{Arbre final. Chaque chemin de la racine vers un \texttt{nil} contient 2 nœuds noirs,
et aucun nœud rouge n'a d'enfant rouge.}
\end{figure}

% ==============================================================================
\section{Résultats}
% ==============================================================================

Machine : Intel Core i5-1135G7, \texttt{g++} 15.2, option \texttt{-O3}. Temps : médianes sur
7 exécutions. Pour 1\,000 et 100\,000 valeurs : les nombres de 1 à $n$ mélangés.

\subsection{Notre arbre (programme 14)}

\begin{table}[H]
\centering
\begin{tabular}{|r|r|r|r|r|}
\hline
$n$ & \textbf{Insertion} & \textbf{Recherche des $n$ valeurs} & \textbf{Hauteur} & $2\log_2(n+1)$ \\
\hline
10      & 1,8 $\mu$s & 0,4 $\mu$s & 4  & 6,9 \\
\hline
1\,000  & 0,17 ms    & 0,05 ms    & 12 & 19,9 \\
\hline
100\,000 & 34,2 ms   & 15,9 ms    & 20 & 33,2 \\
\hline
\end{tabular}
\caption{La hauteur reste toujours sous la borne $2\log_2(n+1)$.}
\end{table}

\subsection{Comparaison avec la STL (programme 15)}

100\,000 clés, trois types de données : \textbf{aléatoires}, \textbf{triées}, et
\textbf{partiellement triées} (triées puis 10\,\% des valeurs échangées au hasard).
On recherche les $n$ clés présentes, puis $n$ clés absentes, et on supprime toutes les clés.
La suppression de notre arbre est vérifiée par 200\,000 opérations au hasard comparées à
\texttt{std::set}.

\begin{table}[H]
\centering
\small
\setlength{\tabcolsep}{3pt}
\begin{tabular}{|l|l|r|r|r|r|r|r|r|}
\hline
\textbf{Données} & \textbf{Structure} & \entete{Insertion} & \entete{Recherche\\clé présente}
  & \entete{Recherche\\clé absente} & \entete{Parcours} & \entete{Suppres-\\sion} & \entete{Total} & \entete{Hau-\\teur} \\
\hline
\multirow{3}{*}{aléatoires} & notre arbre & 45,7 & 18,4 & 23,2 & 3,0 & 43,2 & 133,0 & 20 \\
  & \texttt{map}           & 57,9 & 49,1 & 57,5 & 8,5 & 56,5 & 228,2 & -- \\
  & \texttt{unordered\_map} & \textbf{15,6} & \textbf{1,1} & \textbf{2,4} & \textbf{1,5} & \textbf{9,1} & \textbf{29,6} & -- \\
\hline
\multirow{3}{*}{triées} & notre arbre & 36,5 & 6,4 & 6,9 & 0,9 & 15,7 & 65,9 & 31 \\
  & \texttt{map}           & 32,0 & 17,8 & 17,7 & 1,7 & 13,9 & 83,1 & -- \\
  & \texttt{unordered\_map} & \textbf{8,4} & \textbf{0,7} & \textbf{0,8} & \textbf{0,4} & \textbf{4,9} & \textbf{15,0} & -- \\
\hline
\multirow{3}{*}{partielles} & notre arbre & 32,2 & 11,8 & 13,4 & 1,5 & 22,3 & 81,0 & 29 \\
  & \texttt{map}           & 35,9 & 22,7 & 23,3 & 2,9 & 28,6 & 112,6 & -- \\
  & \texttt{unordered\_map} & \textbf{9,8} & \textbf{0,7} & \textbf{1,3} & \textbf{0,7} & \textbf{6,2} & \textbf{19,0} & -- \\
\hline
\end{tabular}
\caption{Temps en millisecondes pour 100\,000 clés (en gras : le plus rapide).}
\end{table}

\begin{table}[H]
\centering
\footnotesize
\setlength{\tabcolsep}{3pt}
\begin{tabular}{|l|l|r|r|r|r|r|r|}
\hline
\textbf{Données} & \textbf{Structure} & \entete{Comparaisons\\insertion} & \entete{Comparaisons\\recherche}
  & \entete{Comparaisons\\suppression} & \entete{Rota-\\tions} & \entete{Recolo-\\rations} & \entete{Alloca-\\tions} \\
\hline
\multirow{3}{*}{aléatoires} & notre arbre & 3\,202\,671 & 6\,490\,972 & 2\,701\,532 & 100\,163 & 441\,887 & 100\,000 \\
  & \texttt{map}           & 1\,751\,316 & 3\,595\,468 & 2\,428\,631 & -- & -- & 100\,000 \\
  & \texttt{unordered\_map} & 31\,204 & 127\,067 & 100\,000 & -- & -- & 100\,014 \\
\hline
\multirow{3}{*}{triées} & notre arbre & 5\,644\,777 & 6\,538\,194 & 2\,756\,688 & 149\,954 & 799\,734 & 100\,000 \\
  & \texttt{map}           & 2\,972\,387 & 3\,619\,066 & 2\,906\,687 & -- & -- & 100\,000 \\
  & \texttt{unordered\_map} & 13 & 127\,067 & 100\,000 & -- & -- & 100\,014 \\
\hline
\multirow{3}{*}{partielles} & notre arbre & 5\,064\,849 & 6\,825\,470 & 2\,325\,378 & 164\,986 & 641\,413 & 100\,000 \\
  & \texttt{map}           & 2\,682\,423 & 3\,762\,706 & 2\,435\,763 & -- & -- & 100\,000 \\
  & \texttt{unordered\_map} & 8\,492 & 127\,067 & 100\,000 & -- & -- & 100\,014 \\
\hline
\end{tabular}
\caption{Compteurs pour 100\,000 clés. Les rotations et recolorations de la STL se font dans
la bibliothèque : on ne peut pas les compter (--).}
\end{table}

\textbf{Ce qu'on en retient.}
\begin{enumerate}
  \item \textbf{\texttt{unordered\_map} est le plus rapide partout} (4 à 5 fois plus que notre
        arbre au total) : le hachage donne la case directement, en $O(1)$ en moyenne, presque
        sans comparaison. Mais ses clés ne sont pas triées.
  \item \textbf{Notre arbre bat \texttt{std::map} au total} (1,3 à 1,7 fois plus rapide),
        alors que les deux sont des arbres rouge-noir. La différence vient surtout de la
        recherche (1,7 à 2,8 fois plus rapide) : notre nœud ne stocke qu'une clé (pas une paire
        clé/valeur), et notre recherche s'arrête dès qu'elle trouve la clé. Sur les données
        triées, \texttt{map} reste un peu devant pour l'insertion et la suppression.
  \item \textbf{Notre arbre fait pourtant environ 1,8 fois plus de comparaisons} que
        \texttt{map} pour insérer et rechercher : il teste \texttt{==} puis \texttt{<} à
        chaque nœud, \texttt{map} un seul \texttt{<}. Le temps dépend surtout des accès à la
        mémoire, pas du nombre de comparaisons.
  \item \textbf{Les données triées sont plus rapides à parcourir} (les nœuds voisins sont
        utilisés l'un après l'autre), mais elles coûtent plus de rotations et de
        recolorations (150\,000 rotations contre 100\,000), car on insère toujours du même
        côté. La hauteur monte à 31, toujours sous la borne de 33,2.
\end{enumerate}

% ==============================================================================
\section{Conclusion}
% ==============================================================================

L'arbre rouge-noir garde une hauteur en $O(\log n)$ quel que soit l'ordre des données,
grâce à deux outils seulement : la recoloration et la rotation. Même avec des données
triées, pire cas d'un arbre binaire de recherche simple, sa hauteur reste sous
$2\log_2(n+1)$. Face à la STL, notre version écrite à la main est plus rapide que
\texttt{std::map}, mais \texttt{std::unordered\_map} reste bien meilleure quand on n'a pas
besoin de garder les clés triées.

\end{document}
